% ===================================================================
%  Section VI: SP1 Proof Backend
% ===================================================================

The proof backend implements each relation as an SP1 zkVM program following
the host/guest/shared architecture described in \Cref{subsec:sp1}. This
section details the implementation structure and the prove/verify workflow.

\subsection{Backend Organization}

Eight distinct SP1 workspaces reside under \texttt{zk\_backend/}, organized
by relation type:

\begin{itemize}[leftmargin=1.5em]
  \item \texttt{td\_mvp/sp1/}: TD~MVP canonical vector.
  \item \texttt{merkle\_membership/sp1/}: Leaf-hash membership against a
  provenance-bound dataset root, including dataset-size scaling.
  \item \texttt{forward\_td\_mlp/sp1/}: Forward-TD MLP with online argmax and
  target-network value selection.
  \item \texttt{one\_step\_sgd\_tiny/sp1/}: Tiny fixed-point SGD update.
  \item \texttt{short\_trace/sp1/}: Short checkpoint-chain validation.
  \item \texttt{training\_update/sp1/}: Full one-step training update
  (membership + forward + backprop + SGD + checkpoint).
  \item \texttt{training\_fragment/sp1/}: Multi-step fragments for
  $k \in \{1, 4, 8\}$.
  \item \texttt{training\_aggregation/sp1/}: Proof-manifest chunk-chain
  aggregation for $T \in \{32, 64, 128\}$.
\end{itemize}

\noindent Each workspace shares a common pattern: the \emph{shared} crate
defines typed public inputs, private witnesses, fixed-point arithmetic helpers,
SHA-256 hashing wrappers, and the relation-check function. The \emph{guest}
crate reads a typed input via \texttt{sp1\_zkvm::io::read()} and calls the
shared verifier. The \emph{host} crate loads the JSON test vector, constructs
the input, selects execute or prove mode, verifies the proof (if generated),
and writes a compact metrics report.

\subsection{Prove and Verify Workflow}

\Cref{alg:sp1-workflow} summarizes the SP1 prove/verify workflow for a single
relation.

\begin{algorithm}[!t]
\caption{SP1 Relation Prove/Verify Workflow}
\label{alg:sp1-workflow}
\begin{algorithmic}[1]
\Require JSON fixture $F$, mode $\in \{\text{execute}, \text{prove}\}$
\Ensure Verdict $\in \{\text{accept}, \text{reject}\}$, metrics $M$
\State Parse $F$ into public inputs $x$ and witness $w$
\State Compile guest ELF binary
\If{mode = execute}
  \State Run guest natively: $v \gets \mathsf{Guest}(x, w)$
  \State \textbf{assert} $v.\mathsf{ok} = \mathsf{true}$
\ElsIf{mode = prove}
  \State $(\pi, \mathit{pv}) \gets \mathsf{SP1.Prove}(\mathsf{ELF}, x, w)$
  \State \textbf{assert} $\mathsf{SP1.Verify}(\pi, \mathit{pv})$
  \State \textbf{assert} $\mathit{pv}.\mathsf{roots} = x.\mathsf{roots}$
  \Comment{Public-input binding}
  \State Record $M$: cycle count, prove time, verify time, proof size
\EndIf
\State \Return (accept, $M$)
\end{algorithmic}
\end{algorithm}

\subsection{Public-Input Binding}

A critical security property is that the proof public values must match the
public inputs committed before proof generation. For each relation, the host
verifies that roots, hashes, relation identifiers, and claimed outputs in the
proof's public values exactly match those in the original fixture. This prevents
an adversary from generating a valid proof with one set of inputs and
presenting it as proof for a different statement (attack \textbf{A7} from
\Cref{subsec:adversary}).

\subsection{Test Vectors and Provenance}

All SP1 proofs are generated from canonical test vectors stored under
\texttt{zk\_backend/test\_vectors/}. Proof binaries and receipts are not
committed to the repository. Instead, the artifact commits \emph{compact
provenance}: execution metrics, public-input JSON files, witness schemas,
verification reports, proof-artifact policy files, and SHA-256 hash inventories
under \texttt{artifacts/reports/provenance/sp1/}. This design allows reviewers
to verify proof claims without re-running the expensive prove step while
preserving auditability.

\subsection{Tamper Testing in the Backend}

For each SP1 relation, the test suite includes negative (tamper) test cases
that modify specific witness or public-input fields and verify that the guest
rejects the tampered input. The host can optionally disable its own prechecks
(\texttt{--skip-host-precheck}) so that rejection is attributable to the guest
relation logic rather than to host-side JSON filtering. This ensures that
tampering is caught at the correct verification layer (Python oracle, Rust
execute, or SP1 proof verification).
